\section{Related work}
\label{sec:related}

\subsection{Exact multi-objective design-space exploration}
The closest methodological relatives of this work perform exact multi-objective
exploration by reasoning about partial designs. Neubauer et
al.~\cite{neubauer2018exact} encode system-level synthesis in Answer Set Programming
modulo Theories and attach theory propagators that check both constraints and Pareto
dominance on \emph{incomplete} assignments, returning conflict clauses that remove
infeasible and dominated regions early. Their formulation targets heterogeneous
multiprocessors with routing, optimising latency, area and energy, and their pruning is
sound for problems whose objectives are assignment monotonous. In the broader
optimisation literature, multi-objective $A^{*}$ with admissible or
consistent heuristics~\cite{stewart1991moa,mandow2010namoa} provides its theoretical
foundation for bound-based pruning: a partial solution whose optimistic vector is dominated by a complete
solution can be discarded without losing Pareto-optimal completions.

HCA-DSE belongs to this family, and we do not claim safe pruning as such as a
contribution. The difference is the class of objectives. In a distributed CPS the
dominant latency term is contention-dependent: it is a function of tier utilisation,
which no single assignment determines. We therefore construct bounds that are admissible
without assignment monotonicity, by combining a term-wise relaxation with an explicit
decision ordering that makes the contention term exactly computable as soon as its
determinants are fixed (Section~\ref{sec:method}). Avoiding expensive evaluations is also an explicit aim of consistent-approximation
DSE~\cite{neubauer2020consistent}, which combines safe approximations with exact
assessment of promising designs. We therefore do not claim evaluator avoidance
as a new principle. Our study specialises bounds to this CPS model and measures
the evaluation-cost trade-off. The term ``Pareto pruning'' also refers to selecting
a representative subset of an already computed front~\cite{paretopruning2022};
that operation differs from our exact search-space pruning.

\subsection{Design-space exploration for distributed cyber-physical systems}
Herget et al.~\cite{herget2022dse} survey DSE for dCPS and identify the two premises of
this paper: design spaces that are far larger than in classical MPSoC exploration, and
design-point evaluations that are considerably more expensive. CompDSE
\cite{saadatmand2025compdse} uses Y-chart modelling~\cite{kienhuis2002ychart}
and derives workload, platform and mapping models from execution traces of an industrial
lithography machine, evaluating design points with OMNeT++~\cite{varga2008omnet}. It is
an evaluator, not a search strategy: the authors report a manual exploration of three
variations and state that integrating search algorithms is out of scope. The two
approaches are therefore complementary --- CompDSE answers how a design point should be
evaluated, HCA-DSE answers which design points need not be evaluated at all --- and the
per-design-point cost reported there falls in the regime where our crossover analysis
predicts speed-ups of two orders of magnitude.

ContrArc~\cite{xiao2024contrarc} explores CPS architectures with a mixed-integer program
combined with contract refinement checking, pruning through subgraph isomorphism. It
minimises a single objective (cost) subject to requirements rather than constructing a
Pareto front, and its pruning is constraint-based rather than bound-based.

Simulation-based CPS exploration also includes the cyber/control co-design approach of
Amir and Givargis~\cite{amir2020pareto}. Automated workload derivation for industrial
dCPS~\cite{saadatmand2024workload} addresses the input-model construction problem.
Edge CNN exploration~\cite{guo2023} and its robustness extensions~\cite{guo2025}
illustrate application-specific architecture trade-offs at the edge. These studies motivate
broader workloads and reliability models; they do not supply the objective bounds proved here.

\subsection{Exact allocation by integer programming}
\label{sec:related:milp}
A second exact family formulates allocation as an integer program and solves it to
optimality. Kouloumpris et al.~\cite{kouloumpris2024optimization} allocate the tasks of
an application over an edge/hub/cloud triple with binary integer linear programming,
minimising latency or energy under computation, communication, memory and energy
constraints, and validate on a UAV inspection use case with synthetic benchmarks of up to
1000 tasks. Their follow-up~\cite{kouloumpris2026reliability} extends the formulation to
a reliability/latency problem solved exactly and offline for specified objective weights. This certifies a weighted optimum, not enumeration of the entire discrete Pareto front.

These works are exact in a different sense than the present paper, and the distinction
matters for the choice of method. They solve one \emph{allocation} instance over a fixed
infrastructure with a model that is linear in the decision variables, whereas we explore
an \emph{architecture} space in which the number and class of nodes, the link class and
the replication factor are themselves decisions, and in which the latency model contains
a non-linear contention term. A compact mixed-integer linear formulation would require an encoding of the
queueing term. Piecewise-linear approximations can change the objective, but
finite-domain table or disjunctive encodings can in principle retain the exact
values at the cost of model size. Nonlinearity therefore does not rule out an
exact solver baseline. Section~\ref{sec:solver} measures a finite-table SMT baseline, including table construction; a compact symbolic encoding remains outside the comparison.

\subsection{Design-time placement in the edge--cloud continuum}
SPACE4AI-D~\cite{sedghani2024space4aid} is the closest work in terms of the problem
setting: design-time selection of resources and component placement across edge, cloud
and function-as-a-service layers. The problem is formulated as a mixed-integer
non-linear program with a single objective (deployment cost) under response-time
constraints, and solved with randomised greedy, local search, tabu search, simulated
annealing and genetic algorithms; no exactness guarantee is provided. Its validation on
a real deployment reports deviations of $12\%$ in cost and $20\%$ in execution time,
which is a useful reference point for the accuracy reported in
Section~\ref{sec:results}.

Recent Journal of Systems Architecture studies provide complementary context.
Casini et al.~\cite{casini2025edge} combine design-time allocation and timing analysis
with runtime monitoring for real-time edge applications. Workflow scheduling in
MEC-supported CPS~\cite{gene2024} addresses constrained offloading and resource allocation
with a metaheuristic. Dependency-aware microservice offloading has also been studied
on an ICN edge testbed~\cite{ali2026icn}. More recent work examines energy-efficient
cloud--edge inference co-optimisation~\cite{cloud2026coopt}, distributed replica allocation
and load balancing~\cite{filippini2026replica}, and predictive energy-aware continuum
orchestration~\cite{predictive2026}. These are relevant deployment and optimisation
settings, but their application models and dynamic decisions differ from the fixed,
finite architecture space for which we certify an analytical Pareto front.

\subsection{Metaheuristic exploration of edge--cloud deployments}
A large body of work applies evolutionary and swarm metaheuristics to task allocation,
offloading and service placement in edge--cloud systems, typically optimising
combinations of latency, energy and cost. These methods are approximate by construction
and their quality is a function of the evaluation budget. We use NSGA-II
\cite{deb2002nsga2} as implemented in pymoo~\cite{blank2020pymoo} as the budget-matched
representative of this family, and report hypervolume~\cite{zitzler1999hv},
IGD$^{+}$~\cite{ishibuchi2015igdplus} and Pareto recall against the exhaustive ground
truth; the choice of indicators follows the guidance of
Audet et al.~\cite{audet2021indicators}. Our conclusion is not that such methods are inferior in general, but that when
admissible bounds are available exactness is achievable at a small fraction of their
evaluation budget.

Table~\ref{tab:related} summarises the comparison.

\begin{table*}[t]
\centering
\caption{Positioning with respect to the closest prior work. ``Partial-design pruning''
denotes discarding incomplete designs before they are completed and evaluated;
``contention-aware'' denotes objectives that depend on resource utilisation rather than
being additive over decisions. ``Weighted denotes exact optimisation for specified weights, without full-front enumeration.}
\label{tab:related}
\scriptsize
\setlength{\tabcolsep}{3pt}
\begin{tabular}{@{}p{2.5cm}c>{\raggedright\arraybackslash}p{1.95cm}>{\raggedright\arraybackslash}p{2.05cm}cccccc@{}}
\toprule
Work & Year & System class & Design variables & Obj. & Exact & Partial- & Cont.- &
Exp. & Real \\
 & & & & & front & design & aware & eval. & meas. \\
\midrule
Blickle et al.~\cite{blickle1998system} & 1998 & embedded synthesis &
  allocation, binding & 2 & no & no & no & no & no \\
Haubelt \& Teich~\cite{haubelt2003pareto} & 2003 & SoC & component selection & $n$ &
  partly & front comp. & no & no & no \\
Lukasiewycz et al.~\cite{lukasiewycz2008symbolic} & 2008 & embedded synthesis &
  allocation, binding, routing & $n$ & no & SAT-based & no & no & no \\
Neubauer et al.~\cite{neubauer2018exact} & 2018 & heterog.\ MPSoC &
  allocation, binding, routing & 3 & \textbf{yes} & symbolic & no & no & no \\
Herget et al.~\cite{herget2022dse} & 2022 & distributed CPS & survey & --- & --- &
  --- & --- & --- & --- \\
Kouloumpris et al.~\cite{kouloumpris2024optimization} & 2024 & edge/\allowbreak hub/\allowbreak cloud &
  task allocation & 1 & n/a & solver & no & no & UAV case \\
Xiao et al.~\cite{xiao2024contrarc} & 2024 & CPS architectures &
  component selection & 1 & --- & isomorphism & no & no & no \\
Sedghani et al.~\cite{sedghani2024space4aid} & 2024 & computing continuum &
  resources, placement & 1 & no & no & partly & no & \textbf{yes} \\
CompDSE~\cite{saadatmand2025compdse} & 2025 & industrial dCPS &
  cores, hosts, frequency & perf. & --- & none & yes & --- & traces \\
Kouloumpris et al.~\cite{kouloumpris2026reliability} & 2026 & edge/\allowbreak hub/\allowbreak cloud &
  task allocation & 2 & weighted & solver & no & no & case study \\
NSGA-II family~\cite{deb2002nsga2} & --- & edge--cloud & varies & 2--4 & no & no &
  partly & no & rarely \\
\textbf{HCA-DSE (this work)} & 2026 & heterog.\ edge--cloud CPS &
  \textbf{counts, classes, links, placement, replication} & \textbf{4} &
  \textbf{yes} & \textbf{bounds} & \textbf{yes} & \textbf{yes} & \textbf{yes} \\
\bottomrule
\end{tabular}
\end{table*}
